%% kapa: exam=1 round=34 year=2023 period=1 form=A subject=real_estate subject_ko=부동산학원론 seq=8 number=88
%% source: Q-Net 감정평가사 기출문제 (cst00309) artlSeq=5211789 file=raw/1008/5211789/2023년도 제34회 감정평가사 1차시험 1교시.pdf
%% conversion: kordoc 4.24.51 -> tools/gen1.py (auto); review=proofread-sonnet; flags=none
\begin{question}{88}
부동산의 가격탄력성과 균형변화에 관한 설명으로 옳지 않은 것은? (단, 완전탄력적과 완전비탄력적 조건이 없는 경우 수요와 공급법칙에 따르며, 다른 조건은 동일함)
\begin{choices}
  \item 공급이 완전비탄력적일 경우, 수요가 증가하면 균형가격은 상승하고 균형량은 불변이다.
  \item 수요가 완전비탄력적일 경우, 공급이 감소하면 균형가격은 상승하고 균형량은 불변이다.
  \item 수요가 완전탄력적일 경우, 공급이 증가하면 균형가격은 불변이고 균형량은 증가한다.
  \item 공급이 증가하는 경우, 수요의 가격탄력성이 작을수록 균형가격의 하락폭은 크고 균형량의 증가폭은 작다.
  \item 수요가 증가하는 경우, 공급의 가격탄력성이 작을수록 균형가격의 상승폭은 작고 균형량의 증가폭은 크다.
\end{choices}
\end{question}
\answer{5}
